\section{Permiso de bancada con protección de fase seleccionada (GCI-29)}
\label{sec:sd4-permiso-gci29}

Se sustituyen los seis RM35JA32MW de GCI-26 por dos relés SEL-751, uno por grupo, con tres entradas convencionales de corriente nominal 5 A, entradas convencionales de tensión nominal máxima 300 V fase--neutro, fuente de baja tensión 24--48 V DC y tarjeta opcional de cuatro entradas/cuatro salidas estándar, conforme a GCI-32. Se reserva una salida de alarma Form A supervisada por watchdog y una salida estándar Form A para permiso de bancada. Es una configuración publicada del SEL-751; la referencia de pedido configurada y el plano de bornes se incorporan al expediente de compra antes de energizar. No se adquiere la variante LEA ni se usa como referencia de pedido el ejemplo de figura que contiene entradas LEA. La alimentación sigue siendo independiente por camino desde las dos QUINT de CEN-26. Se añaden dos Phoenix PTCB E1 24DC/2A NO, referencia 2909903, exclusivamente para los SEL; se conservan los cuatro PTCB de 1 A para XPS. El SEL requiere menos de 30 W DC, de modo que 1 A a 24 V no cubre su cota y no se reutiliza esa protección. El PTCB de 2 A admite 18--30 V DC y disipa menos de 0,8 W nominales; no limita recarga \parencites[pp.~2, 30--35 y 38--40]{lote29-sel751}{lote29-ptcb2}.

Se conservan los seis TC METSECT5CC015 de relación 150/5 A y clase 0,5/3 VA de GCI-26. Con 4 m de lazo de cobre 1,5 mm$^2$, factor térmico de proyecto 1,14 y reserva de bornes 0,30 VA, la carga secundaria es como máximo $1,1667(1,14)+0,10+0,30=1,7300$ VA: la entrada SEL publica menos de 0,10 VA a 5 A. Se retira la carga de entrada RM35 de 0,375 VA; no se suman ambas. Cada secundario dispone de bloque cortocircuitable de pruebas, sin fusible ni maniobra que lo abra bajo corriente primaria; se verifican polaridad, continuidad, carga instalada y clase efectiva. Se seleccionan bornes seccionables/cortocircuitables de función de prueba como partida de instalación, cuya referencia de compra queda por individualizar. La toma de tensión se realiza en estrella sobre las tres fases y neutro del tablero nuevo 400/230 V, con derivación protegida por fase, segregación y pruebas de ausencia/inversión de tensión. La referencia/poder de corte de esas protecciones y la detección efectiva de fallos de TC quedan por desarrollar. La alarma interna no acredita por sí sola integridad de TC, bornes o bobinas \parencites[p.~33]{lote29-sel751}[pp.~1--2]{lote29-selhealth}[p.~1]{lote26-ct}.

Se fija el elemento de sobrecorriente de fase 50P a 3,75 A secundarios, equivalente nominal a 112,5 A primarios, con retardo intencional 0,00 s. Se fija sobretensión 59P a 239,00 V fase--neutro, también con retardo 0,00 s. Son ajustes propios dentro de rangos publicados, no valores de fábrica. SEL publica para 50P error de $\pm3\,\%$ del ajuste más $\pm0,02 I_{\rm NOM}$ en régimen y $\pm5\,\%$ más $\pm0,02 I_{\rm NOM}$ en transitorio. Con $I_{\rm NOM}=5$ A y error de relación TC de 0,5\,\%, los extremos superiores de disparo tardío resultan
\[
 I_{\rm reg}=\frac{(3,75\cdot1,03+0,10)30}{0,995}=119,4724~{\rm A},
 \qquad I_{\rm tr}=\frac{(3,75\cdot1,05+0,10)30}{0,995}=121,7337~{\rm A}.
\]
El error publicado de tensión $\pm1\,\%$ del ajuste más $\pm0,5$ V produce extremo superior $V_{\rm f}=239(1,01)+0,5=241,89$ V. En la envolvente conjunta, después de responder los elementos y con los tres circuitos de medida íntegros, $\sum V_{\rm f} I_{\rm f}\le3(241,89)(121,7337)/1000=88,3385$ kVA, inferior a 90 kVA; $P\le S$ también acota kW. Esta desigualdad admite desequilibrio sólo si cada fase satisface sus cotas; no se obtiene usando únicamente $\sqrt3VI$ con una fase representativa. No cubre saturación de TC, armónicos fuera del alcance del filtro de protección, incertidumbres no publicadas de instalación ni la ventana transitoria previa a retirada \parencite[pp.~38--40]{lote29-sel751}.

El permiso se da por una conjunción local: alarma Form A sana, protección habilitada y medidas comprobadas, ausencia de 50P/59P y ausencia de causa ambiental o parada. La alarma Form A y la salida de permiso se ponen en serie entre PTCB de 1 A y A1 de SR-B; A2/B2 conserva retorno propio. Form A abre al perder alimentación o al detectar fallo interno/watchdog. El permiso abre ante disparo y queda enclavado hasta investigación y rearme manual; recuperación de tensión o comunicación no autoriza una bancada. La selección de contactos, bornes y plantilla se actualiza en GCI-32. La salida se mantiene abierta durante puesta en servicio y arranque: la ficha indica arranque aproximado de 5--10 s, que no es máximo garantizado ni temporización permisiva. El rearme exige alarma sana y validación local de protección, sin esperar sólo diez segundos ni usar el LED ENABLED como sustituto del watchdog. Los nombres de variables, bornes y lógica ejecutable se fijan con el manual vigente de la configuración y se prueban antes de habilitar; el contrato lógico de este párrafo no se presenta como fichero de ajustes ya cargado \parencites[pp.~34 y 38]{lote29-sel751}[pp.~1--2]{lote29-selhealth}.

La retirada de SR-B conserva los dos LC1D32P7 en serie, sus espejos NC y rearme de CEN-26. No retira SR-R por un exceso atribuible únicamente a bancada. Si el servicio sin bancada ya excede la capacidad del grupo, abrir bancada no resuelve el exceso: se impide solicitarla y se mantiene pendiente la coordinación del servicio, recarga y climatización. La limitación manual de diez pasos sigue vigente; el SEL no convierte LC10/VCS en controlador automático de potencia ni resuelve consumo a baja carga y reposición de combustible.

GCI-32 sustituye la interfaz híbrida por contacto estándar y retira la suma anterior de 184 ms. El tiempo publicado de protección con salida híbrida no acredita el recorrido con contacto estándar. Desde pérdida efectiva de 24 V en A1 de XPS sólo se conservan 139 ms parciales de XPS/contactor; elemento, lógica, salida y extinción requieren presupuesto completo y prueba integral. No se aplica fuga de salida híbrida ni su tiempo de apertura al contacto estándar. La cota de 88,3385 kVA anterior sigue siendo posterior a respuesta: no autoriza la bancada durante una ventana transitoria sin comprobar. La aceptación y el permiso retirado siguen GCI-32 \parencites{lote32-selmanual2024}{lote29-sel751}{lote26-xpsmanual}{lote26-contactor}.

La recepción prevista incluye inyección secundaria de las tres fases, error de relación/carga de TC, cada umbral por separado y conjuntamente, fase sin tensión, pérdida de alimentación, alarma watchdog, arranque, rearme y fallo de contacto final. El registro mide corriente/tensión reales y tiempo hasta retirada de potencia, incluidas transición del grupo y escalón manual. La aceptación debe demostrar ausencia de exceso sostenido y tolerancia del generador durante la ventana; no se libera el límite de 90 kVA por la suma anterior. No se atribuye redundancia de protección completa ni nivel SIL al conjunto por disponer de dos relés en caminos distintos.

\section{Demanda de control y condición de recarga (ECI-29)}
\label{sec:sd4-control-eci29}

La sustitución retira 3,6 W DC de seis RM35 y añade 60 W de cota de dos SEL y 1,6 W de dos PTCB: $14,8-3,6+60+1,6=72,8$ W DC, 36,4 W por QUINT antes de sensores ambientales. Cada fuente nominal de 240 W cubre esa carga estática; arranque, selectividad y autonomía se comprueban con los receptores restantes. No se agrega de nuevo este DC al AC: se conserva reserva de entrada nominal de fuentes 548 VA entre ambos caminos, y totales de 558 VA con bancada retirada/596 VA con ambas permitidas de CEN-26. Los sensores y actuadores adicionales deben incorporarse por camino/estado antes del balance definitivo. No se atribuye eficiencia nominal de QUINT a su consumo parcial \parencites[p.~34]{lote29-sel751}{lote29-ptcb2}[pp.~3--4]{p6-puestos-quint}.

El ajuste requerido de 20 A agregados del 93T sigue sin perfil inferior/resolución autorizado, ni supervisión efectiva por cadena. La cota ideal de gas a 60 A y 900 m$^3$/h supera el umbral crítico de proyecto; disponer de ese caudal no habilita recarga. Se aplica GAS-28 y se conserva el enclavamiento ambiental independiente. El SEL seleccionado protege demanda de grupo al retirar bancada: no controla corriente DC de batería ni habilita por sí solo recarga, y retirar rectificador carga la batería. Se mantiene el permiso de recarga retirado hasta seleccionar perfil OEM y desarrollar supervisión efectiva; la obligación de descarga--recarga--nuevo corte sigue en CEN/ECI \parencite[p.~98]{p7-eaton-93t-guia}.

\section{Cribado integral de sustitución UPS (CEN-29)}
\label{sec:sd4-cribado-cen29}

Se analiza como candidato, sin incorporarlo al inventario vigente, el Eaton 93PS de 40 kW/40 kVA con dos UPM de 20 kW, sin MBS ni baterías internas, referencia publicada BD04A0206A01000000. Tres marcos por camino darían 120 kW/120 kVA; seis marcos medirían 480$\times$750$\times$1750 mm cada uno y sumarían 1248 kg de masa de catálogo. La especificación revisión 016 publica recarga configurable 2--50 A por marco hasta 80\,\% de carga y 2--30 A por encima de 80\,\%, con valor por defecto 10 A. Se usa sólo el extremo publicado 2 A para cribado: tres marcos sumarían 6 A por camino; el valor por defecto sumaría 30 A, superior al cribado sin margen de 28,4293 A de GAS-28. No se inventa resolución de ajuste ni se declara configurado el cargador \parencites[pp.~2 y 5]{lote29-93psspec}{lote29-93pssku}.

Tres cadenas de 40 CSB HRL12540W por camino requerirían 240 bloques y seis armarios. La cota precedente de autonomía por bloque a 60 min, con factores 0,85/0,90, sería $120(1129)(0,85)(0,90)/1000=103,6422$ kW por camino, inferior a 120 kW nominales. Cada cadena pesa $40(43,9)=1756$ kg; seis suman 10536 kg sin armarios. Con reserva propia de 300 kg/armario y apoyo de 1,6 m$^2$, una cadena ocuparía 2056 kg y requeriría $2056(9,80665)/1,6/1000=12,6015$ kPa de comprobación estructural; la reserva no es masa de fabricante del armario. La tercera cadena por camino exige ampliar el trazado y comprobar suelo, servicio y evacuación. No se añade sólo a la tabla para declarar autonomía \parencite{csbHRL12540WRA260131}.

La entrada máxima publicada de 76 A por marco suma 228 A por camino a 400 V, aproximadamente 157,9630 kVA, y no es compatible con el grupo de 90 kVA como máxima simultánea. Deben quedar seleccionados y verificados reparto/limitación de entrada, perfil de carga, bypass, cargas de entrada única y respuesta del grupo antes de adoptar este candidato. La eficiencia publicada del marco de 40 kW es 95,7/96,1/96,2/95,5\,\% a 100/75/50/25\,\% de carga, respectivamente; no se extrapola por debajo de 25\,\% ni se deduce calor de recarga no tabulado. El crecimiento ya necesario por recalentamiento y ventilación, junto con frío y bombas pendientes, impide seleccionar una UPS mayor desde un subtotal. CEN-29 deja desarrolladas dimensiones, masa, autonomía y causas cuantificadas de rechazo provisional; la UPS 93T vigente no obtiene conformidad a plena carga por esta alternativa \parencite[pp.~2--5]{lote29-93psspec}.

